\section{関連研究}\label{関連研究}
\begin{comment}
・自己開示をシステムが促すことで親密度が向上するんだよーっていうKDDIの研究をかく
・カラオケでの楽曲推薦とかも必要そうなら挙げる
・カラオケにはシステムがあるから，介入させるのが自然だよね（既存のDAMやJOYSOUNDでrecommendされてることを記す）

→ システムを使うことが有効そう だと示唆する
\end{comment}

\subsection{自己開示の促進支援に関する研究}
池田ら\cite{自己開示の促しによるコミュニケーション支援システム}の研究は，自己開示を促進するコミュニケーション支援システムの有効性を示した重要な知見を提供している．池田らは婚活パーティーを想定し，初対面の2人の参加者間でタブレット端末に話題を表示するシステムを開発した．話題には「趣味」や「休日の過ごし方」といった自己開示項目を用い，(1)自己開示の深さを考慮して段階的に話題を提示する手法，(2)ランダムに話題を提示する手法，(3)システムを用いない条件の3つの方法で12分間の会話実験を行った．192名の参加者による大規模な実験の結果，システムによる自己開示の促進が，参加者間の親密度向上に有意な効果をもたらすことを実証した．また，段階的な提示手法では，提示された話題への興味や言及のしやすさ，システムの継続利用意向といった評価指標において，ランダム提示手法より高い評価を得ることが確認された．ただし，システムへの過度な依存や自然なコミュニケーションの阻害といった課題も指摘されており，ユーザの自由度を保ちながら必要な時にのみ支援を提供することの重要性が示唆されている．

\subsection{カラオケ環境における支援システム}
カラオケ環境は，システムによる介入が自然に受け入れられる特徴を持つ．選曲の際にタブレット端末を使用することが一般的であり，システムからの情報提示に対する抵抗感が低いと考えられる．実際に，田中ら\cite{tanaka_2020}によって開発されたレコメンドシステムがDAM\cite{DAM}のカラオケシステム「LIVE DAM Ai」に導入されているように，楽曲推薦による支援が実用化されている．このような特性は，自己開示を支援するシステムを導入する上で大きな利点となり得る．

\vskip.5\baselineskip

本研究では，これらの知見を踏まえ，カラオケ環境の特性を活かしながら，自然なコミュニケーションを阻害しない形での自己開示支援システムの設計を目指す．